\section{Определение полуторалинейной формы}

\begin{definition}[Полуторалинейная форма]
Пусть $V$ --- линейное пространство над $\mathbb{C}$ ($F = \mathbb{C}$).

Функция $g(x, y): V \times V \to \mathbb{C}$ называется \textbf{полуторалинейной формой}, если выполнены 4 свойства:
\begin{enumerate}
    \item $g(x_1 + x_2, y) = g(x_1, y) + g(x_2, y),\quad \forall x_1, x_2, y \in V$;
    \item $g(\lambda x, y) = \lambda g(x, y),\quad \forall x, y \in V,\ \forall \lambda \in \mathbb{C}$;
    \item $g(x, y_1 + y_2) = g(x, y_1) + g(x, y_2),\quad \forall x, y_1, y_2 \in V$;
    \item $g(x, \mu y) = \bar{\mu}\, g(x, y),\quad \forall x, y \in V,\ \forall \mu \in \mathbb{C}$.
\end{enumerate}

Форма называется \textbf{эрмитово симметричной}, если $g(x, y) = \overline{g(y, x)}$.
\end{definition}
